% Chapter — Gravity-gradient torque (FR-1). Follows the mandatory FR-29
% six-section template: purpose; assumptions and domain of validity;
% derivation; discretization and implementation; validation evidence;
% references. The equation labels eq:gravgrad:torque, eq:gravgrad:pitch,
% eq:gravgrad:libfreq, and eq:gravgrad:pendulum are echoed verbatim in
% comments in cpp/src/models/gravgrad.cpp and cpp/tests/test_gravgrad.cpp
% (FR-29 equation-label-to-code traceability); renaming them breaks that
% trace.
\chapter{Gravity-Gradient Torque}
\label{ch:gravgrad}

\section{Purpose}
\label{sec:gravgrad:purpose}

This chapter documents the gravity-gradient torque model (FR-1): the torque
a finite rigid body experiences about its center of mass because the
central body's gravitational field is stronger on the near side than on the
far side. It is the one environmental torque the PRD models --- on-orbit
aerodynamic, solar-radiation-pressure, and magnetic disturbance torques are
Non-Goals (their exclusion is bounded in
Section~\ref{sec:gravgrad:domain}). The model is implemented in
\texttt{cpp/src/models/gravgrad.cpp} with its interface in
\texttt{cpp/include/star/models/gravgrad.hpp}; it is purely algebraic in
its inputs $(\mu, \vv{r}^{I}, \quat{q}_{i2b}, \mat{I})$ --- the
gravitational parameter of the current central body, the
central-body-to-vehicle position in the inertial frame, the attitude, and
the inertia tensor, all quantities the FR-1 state and mass-property layers
already carry --- and depends on no environment or vehicle structure. The
chapter also derives the pitch-libration frequency of a gravity-gradient
stabilized body, the analytic reference for Phase~4 exit criterion~9.

\section{Assumptions and domain of validity}
\label{sec:gravgrad:domain}

Assumptions:
\begin{enumerate}
  \item \textbf{Point-mass central field.} The torque is computed from the
    spherically symmetric $-\mu\vv{r}/r^3$ field. The field's harmonic
    structure (Chapter~\ref{ch:gravity}) also has a gradient across the
    body, but its contribution is smaller than the retained term by the
    relative size of the harmonic field itself,
    $\mathcal{O}\bigl(J_2 (R_\oplus/r)^2\bigr) \sim 10^{-3}$ in LEO, and
    is neglected --- consistent with the standard formulation.
  \item \textbf{First-order tidal expansion.} The field is linearized
    across the body: for a vehicle of extent $c$ at distance $r$, the
    first neglected term is smaller by $\mathcal{O}(c/r)$
    ($\sim 10^{-6}$ for $c \sim \qty{10}{\metre}$ in LEO). The expansion
    is the same tidal limit as the third-body discussion of
    Chapter~\ref{ch:thirdbody}, applied here to the attitude channel.
  \item \textbf{Torque about the instantaneous center of mass}, with
    $\mat{I}$ the same center-of-mass, body-axes inertia tensor the
    dynamics of Chapter~\ref{ch:rigidbody} uses, and $\vv{r}$ the vector
    \emph{from the current central body to the vehicle} (FR-1: "about the
    current central body, using states and inertia already available").
    Only the direction of $\vv{r}$ and its magnitude enter.
  \item \textbf{Unit attitude quaternion.} $\quat{q}_{i2b}$ carries the
    caller's unit-norm invariant (maintained by the FR-1 post-step
    normalization); a non-unit quaternion scales the returned torque by
    $\norm{\quat{q}}^{4}$ (the DCM enters twice), a relative error of
    $4(\norm{\quat{q}} - 1)$ for small deviations --- negligible under the
    drift bound of Chapter~\ref{ch:rigidbody} but stated so the invariant
    is explicit.
  \item \textbf{Prescribed-orbit libration analysis.} The libration
    derivation of Section~\ref{sec:gravgrad:libration} and its validation
    test place the body on a \emph{prescribed} circular orbit of radius
    $r$ with mean motion $n = \sqrt{\mu/r^3}$: attitude does not feed back
    into translation. This is the module-level validation configuration;
    the fully coupled translational--rotational vehicle lands with the
    Phase~4 integration workstream, and the attitude-to-orbit coupling
    neglected here is of relative order $(c/r)^2$ in the orbit.
\end{enumerate}

Excluded-torque bounds, recorded once so the exclusion is quantitative:
for a vehicle of characteristic area $A$, center-of-pressure offset $d$,
and residual dipole $\vv{m}$, the excluded torques are bounded by
$\tau_{\mathrm{srp}} \le P_{1\mathrm{au}} A d\,(1 + C_R)$ with
$P_{1\mathrm{au}} = \qty{4.54e-6}{\newton\per\metre\squared}$
(Chapter~\ref{ch:srp}), $\tau_{\mathrm{aero}} \le \tfrac{1}{2}\rho v^2
C_D A d$ with $\rho$ from the atmosphere chapters
(Chapters~\ref{ch:ussa76}--\ref{ch:harrispriester}), and
$\tau_{\mathrm{mag}} = \norm{\vv{m} \times \vv{B}}$, parametric in the
vehicle's residual dipole and the local field (no field model ships with
the core). For a representative $\qty{100}{\kilogram}$-class LEO vehicle
($A \sim \qty{1}{\metre\squared}$, $d \sim \qty{0.1}{\metre}$,
$\mat{I} \sim \qty{100}{\kilogram\metre\squared}$) the modeled
gravity-gradient torque at $r = \qty{7e6}{\metre}$ is of order
$3\mu\,\Delta I/(2 r^3) \sim \qty{7e-5}{\newton\metre}$, one to two orders
above the SRP bound ($\sim \qty{1e-6}{\newton\metre}$) and above the aero
bound outside the densest drag regimes.

Domain of validity: any state with $\vv{r} \neq \vv{0}$. Out-of-domain
behavior, matching the code exactly: at $\vv{r} = \vv{0}$ the
normalization divides by zero and IEEE non-finite values propagate; the
model does not throw (deterministic-loop rule; a spacecraft coincident
with a body center is excluded by configuration validation, as in
Chapter~\ref{ch:thirdbody}).

\section{Derivation}
\label{sec:gravgrad:derivation}

\subsection{The torque integral}

Let $\vv{x}$ run over the body's mass elements measured from the center of
mass ($\int \vv{x}\,\mathrm{d}m = \vv{0}$), and let $\vv{r}$ point from the
central body to the center of mass, $r = \norm{\vv{r}}$, $\hat{\vv{r}} =
\vv{r}/r$. Each element feels the point-mass field of
equation~\eqref{eq:twobody:newton}, so the torque about the center of mass
is
\begin{equation}
  \vv{\tau} = \int \vv{x} \times
    \left( -\,\mu\, \frac{\vv{r} + \vv{x}}{\norm{\vv{r} + \vv{x}}^{3}}
    \right) \mathrm{d}m
  = -\,\mu \int \frac{\vv{x} \times \vv{r}}
      {\norm{\vv{r} + \vv{x}}^{3}}\, \mathrm{d}m ,
  \label{eq:gravgrad:torqueint}
\end{equation}
using $\vv{x} \times \vv{x} = \vv{0}$. Expanding the denominator to first
order in $\norm{\vv{x}}/r$,
\begin{equation}
  \norm{\vv{r} + \vv{x}}^{-3}
    = r^{-3} \Bigl( 1 - 3\, \frac{\hat{\vv{r}} \cdot \vv{x}}{r}
      + \mathcal{O}\bigl(\norm{\vv{x}}^2/r^2\bigr) \Bigr),
  \label{eq:gravgrad:expansion}
\end{equation}
the zeroth-order term integrates to $(\int \vv{x}\,\mathrm{d}m) \times
\vv{r} = \vv{0}$ (center-of-mass property), leaving
\begin{equation}
  \vv{\tau} = \frac{3\mu}{r^{3}} \int (\hat{\vv{r}} \cdot \vv{x})
    \,(\vv{x} \times \hat{\vv{r}})\, \mathrm{d}m
    + \mathcal{O}\!\left(\frac{c}{r}\right)\vv{\tau} .
  \label{eq:gravgrad:firstorder}
\end{equation}
The integral is an inertia-tensor contraction: with $\mat{I} = \int
(\norm{\vv{x}}^2 \mat{I}_3 - \vv{x}\vv{x}^{\mathsf T})\,\mathrm{d}m$,
\begin{equation}
  \hat{\vv{r}} \times (\mat{I}\hat{\vv{r}})
  = \int \Bigl( \norm{\vv{x}}^2\, \hat{\vv{r}} \times \hat{\vv{r}}
    - (\vv{x} \cdot \hat{\vv{r}})\, (\hat{\vv{r}} \times \vv{x})
    \Bigr) \mathrm{d}m
  = \int (\hat{\vv{r}} \cdot \vv{x})\,(\vv{x} \times \hat{\vv{r}})\,
    \mathrm{d}m ,
  \label{eq:gravgrad:identity}
\end{equation}
so the gravity-gradient torque, evaluated with every vector resolved in
body axes ($\hat{\vv{r}}^{B} = \dcm{I}{B}\,\hat{\vv{r}}^{I}$, since
$\mat{I}$ lives in body axes), is
\begin{equation}
  \boxed{\;
  \vv{\tau}^{B}_{gg}
    = \frac{3\mu}{r^{3}}\;
      \hat{\vv{r}}^{B} \times \bigl( \mat{I}\, \hat{\vv{r}}^{B} \bigr) ,
  \;}
  \label{eq:gravgrad:torque}
\end{equation}
the standard result (Markley and Crassidis~\cite{markley2014}). Three
structural properties follow directly and are gated bit-exactly by the
property tests: the torque vanishes when $\mat{I}$ is spherical
($\mat{I}\hat{\vv{r}} \parallel \hat{\vv{r}}$) and when $\hat{\vv{r}}$ is
any principal axis (eigenvector) of $\mat{I}$; and its magnitude is
bounded by $\norm{\vv{\tau}} \le \tfrac{3\mu}{2r^{3}}
(I_{\max} - I_{\min})$, the extremum occurring with $\hat{\vv{r}}$ at
$45^{\circ}$ between the extreme principal axes.

\subsection{Pitch libration on a circular orbit}
\label{sec:gravgrad:libration}

Place the body on a prescribed circular orbit of radius $r$ in the
inertial $xy$ plane, orbit angle $\nu = n t$, $n = \sqrt{\mu/r^3}$, and
let the body track the local-vertical frame with a pitch offset: body
$\hat{\vv{y}}_B$ along the inertial orbit normal $+\hat{\vv{z}}_I$ (the
pitch axis), body $\hat{\vv{z}}_B$ at angle $\theta$ from the zenith
direction $\hat{\vv{r}}$, so the body's total rotation angle about the
orbit normal is $\alpha = \nu + \theta$ and $\vv{\omega}^{B} = (0,\,
n + \dot{\theta},\, 0)$. Resolving the zenith direction in body axes
gives $\hat{\vv{r}}^{B} = (-\sin\theta,\, 0,\, \cos\theta)$, and for a
principal-axes inertia $\mat{I} = \mathrm{diag}(I_x, I_y, I_z)$ ($x$
along-track, $y$ pitch, $z$ nadir at $\theta = 0$),
equation~\eqref{eq:gravgrad:torque} reduces to a pure pitch torque
$\vv{\tau}^{B} = (0, \tau_y, 0)$ with $\tau_y = -3 n^2 (I_x - I_z)
\sin\theta \cos\theta$. This planar configuration is \emph{exactly}
self-consistent: with $\vv{\omega}$ on the pitch axis the gyroscopic term
$\vv{\omega} \times (\mat{I}\vv{\omega})$ vanishes identically, the
torque stays on the pitch axis, and roll/yaw remain exactly zero.
Euler's equation~\eqref{eq:rigidbody:euler} about the pitch axis, with
$\ddot{\alpha} = \ddot{\theta}$ on the circular orbit, is then the exact
pendulum
\begin{equation}
  I_y\, \ddot{\theta}
    = -\,3 n^2 (I_x - I_z) \sin\theta \cos\theta
    = -\,\tfrac{3}{2}\, n^2 (I_x - I_z) \sin 2\theta .
  \label{eq:gravgrad:pitch}
\end{equation}
Linearizing for small $\theta$ gives simple harmonic pitch libration at
\begin{equation}
  \boxed{\;
  \omega_{\mathrm{lib}}
    = n\, \sqrt{\frac{3\,(I_x - I_z)}{I_y}} ,
  \;}
  \label{eq:gravgrad:libfreq}
\end{equation}
stable (oscillatory) precisely when $I_x > I_z$ --- the along-track moment
exceeding the nadir moment, i.e.\ the long axis pointing at the central
body, the classical gravity-gradient-stable configuration. This is the
analytic value of Phase~4 exit criterion~9.

At finite amplitude the measured frequency is predictably lower. In
$u = 2\theta$, equation~\eqref{eq:gravgrad:pitch} is $\ddot{u} =
-\omega_{\mathrm{lib}}^2 \sin u$; its energy integral $\tfrac{1}{2}
\dot{u}^2 = \omega_{\mathrm{lib}}^2 (\cos u - \cos u_0)$ for release from
rest at amplitude $u_0$ gives the quarter period $\int_0^{u_0}
\mathrm{d}u / \dot{u}$, which the substitution $\sin(u/2) =
\sin(u_0/2)\sin\chi$ reduces to the complete elliptic integral of the
first kind $K(m) = \int_0^{\pi/2} (1 - m \sin^2\chi)^{-1/2}\,
\mathrm{d}\chi$ with parameter $m = \sin^2(u_0/2)$. The oscillation
frequency at initial pitch offset $\theta_0$ (so $u_0 = 2\theta_0$, $m =
\sin^2\theta_0$) is therefore
\begin{equation}
  \omega_{\mathrm{pend}}(\theta_0)
    = \omega_{\mathrm{lib}}\, \frac{\pi}{2\,K(\sin^2\theta_0)}
    = \omega_{\mathrm{lib}}
      \Bigl(1 - \tfrac{1}{4}\theta_0^{2}
      + \mathcal{O}(\theta_0^{4})\Bigr) ,
  \label{eq:gravgrad:pendulum}
\end{equation}
a $-2.5 \times 10^{-5}$ relative offset at the committed $\theta_0 =
\qty{0.01}{\radian}$ --- forty times inside the exit-criterion gate, and
the basis of the tighter secondary check in the validation test.

\section{Discretization and implementation}
\label{sec:gravgrad:implementation}

There is no discretization: the model is a pointwise algebraic map from
$(\mu, \vv{r}^{I}, \quat{q}_{i2b}, \mat{I})$ to $\vv{\tau}^{B}$.
Implementation notes:
\begin{enumerate}
  \item \textbf{Module and traceability.} The model lives in
    \texttt{cpp/src/models/gravgrad.cpp}; the source comments echo
    \texttt{eq:gravgrad:torque} verbatim, and the validation test echoes
    \texttt{eq:gravgrad:pitch}, \texttt{eq:gravgrad:libfreq}, and
    \texttt{eq:gravgrad:pendulum}.
  \item \textbf{Evaluation order.} The position is normalized in the
    inertial frame ($\hat{\vv{r}}^{I} = \vv{r}/r$) and rotated with
    \texttt{star::rotation::quat\_transform}, the single project code path
    for frame transformation (Chapter~\ref{ch:rotations}), then contracted
    per equation~\eqref{eq:gravgrad:torque}. Fixed order, IEEE basic
    operations plus one square root, no allocation (D-10).
  \item \textbf{Scale factor once.} $3\mu/r^{3}$ is formed as one scalar
    and applied to the cross product, so the structural zeros of the
    property tests (spherical $\mat{I}$, principal-axis alignment) are
    exact: every floating-point product in the cross-product difference
    either carries a zero factor or subtracts two identically rounded
    copies of the same real product.
  \item \textbf{Data ownership.} $\mu$ is supplied by the caller (the
    force-composition layer passes the current central body's value); the
    module defines no body constants.
\end{enumerate}

\section{Validation evidence}
\label{sec:gravgrad:validation}

Table~\ref{tab:gravgrad:validation} lists the tests that constitute this
chapter's validation evidence. Each test identifier is machine-checked to
exist in the test tree by \texttt{scripts/lint\_chapters.py}; golden
provenance and tolerance derivations are recorded in
\texttt{tests/golden/attitude/manifest.toml}.

\begin{table}[htbp]
  \centering
  \caption{Validation evidence for the gravity-gradient model (doctest,
    \texttt{cpp/tests/test\_gravgrad.cpp}).}
  \label{tab:gravgrad:validation}
  \begin{tabular}{lp{8.6cm}}
    \toprule
    Test ID & Acceptance basis \\
    \midrule
    \testid{gravgrad_torque_golden} &
      Equation~\eqref{eq:gravgrad:torque} against extended-precision
      references (mpmath, 60 significant digits) at six committed states
      --- generic LEO/lunar/GEO geometries with full and diagonal inertia
      tensors, plus the planar pitch case cross-checked at generation
      time against the closed form of
      equation~\eqref{eq:gravgrad:pitch} --- to norm-relative
      $\le 10^{-12}$. \\
    \testid{gravgrad_exact_zero_geometries} &
      Structural zeros bit-exact: spherical inertia (any attitude and
      position) and central-body direction along a principal axis both
      return exactly $\vv{0}$, per the properties of
      equation~\eqref{eq:gravgrad:torque}. \\
    \testid{gravgrad_libration_frequency} &
      Phase~4 exit criterion 9: the libration frequency of a
      gravity-gradient-stabilized body on the prescribed circular
      reference orbit, measured from event-located pitch zero crossings
      over $\sim$11 periods, matches the analytic
      $\omega_{\mathrm{lib}}$ of equation~\eqref{eq:gravgrad:libfreq}
      to within 0.1\,\% (the known finite-amplitude offset is
      $2.5 \times 10^{-5}$), and matches the amplitude-corrected
      $\omega_{\mathrm{pend}}$ of equation~\eqref{eq:gravgrad:pendulum}
      to within $10^{-7}$ relative. \\
    \bottomrule
  \end{tabular}
\end{table}

\section{References}
\label{sec:gravgrad:references}

This chapter relies on Markley and Crassidis~\cite{markley2014} for the
standard gravity-gradient torque result (re-derived above from the tidal
expansion) and on the IERS Conventions (2010)~\cite{iers2010} for the
geocentric gravitational constant used by the committed validation
geometry. Full bibliographic details appear in the bibliography.
